\chapter{Összefoglalás}

\section{Megvalósított funkciók}
Végeredményben a rendszer egy komplex kvízalkalmazást valósít meg, mely egyesíti a többjátékos kompetitív élményt oktatási funkciókkal. A rendszer központi elemei a valós idejű rangsorolt játékmenet, a fejlődés nyomonkövetését elősegítő rendszerek valamint az adminisztrációs rendszer.

A rendszer a következő funkciókat valósítja meg:
\begin{itemize}
    \item Reszponzív viselkedésforma
    \item Regisztráció, Autentikáció
    \item Játékosprofilok
    \item ELO rangrendszer
    \item XP rendszer
    \item Daily streak (napi sorozat) számítás
    \item Winstreak (nyerési sorozat) számítás
    \item Egyjátékos gyakorló meccsek
    \item Többjátékos meccskeresési funkció (SBMM)
    \item 1 vs. 1 rangsorolt mérkőzések
    \item Játékos -és kérdés bejelentési funkciók
    \item Jogosultságrendszer
    \item Adminisztrációs rendszer
\end{itemize}

\section{Nehézségek és tanulságok}
A rendszer komplexitása miatt több helyen is szükség volt kreativitásra a problémák megoldása érdekében:

\begin{itemize}
    \item \textbf{Egyszerű csalásra használható scriptek kiküszöbölése:} A lehető legminimálisabb információkat szabad közölni a klienssel, annak érdekében, hogy az automatizált csalásokat - amelyek a hálózati kéréseket figyelik - kiküszöbölje a rendszer.
    \item \textbf{Valós idejű játékmenet:} Edge case-ek kiküszöbölése: játékos nem fogadja el a meccset, kilép, visszacsatlakozik, több lapon használja az alkalmazást. 
    \item \textbf{Meccskeresés:} Fair, jól balance-olt rendszer kialakítása, minimális erőforrásigénnyel.
    \item \textbf{Tesztelés:} A sok funkció manuális tesztelése nehézkes volt, ami rávilágított az automatizált tesztek fontosságára.
\end{itemize}

\newpage

\section{Továbbfejlesztési lehetőségek}
A rendszer jelenlegi állapota egy MVP-t (Minimum viable product) valósít meg. A rendszer számos továbbfejlesztési lehetőséggel rendelkezik, amelynél a fő szempontok a következők: csalás kiküszöbölése, SBMM továbbfejlesztése, játékmódok bővítése.

Néhány példa továbbfejlesztési lehetőségekre:
\begin{itemize}
    \item \textbf{Csalás kiküszöbölése:} Rangsorolt meccs folyamán több adat rögzítése, annak érdekében, hogy könnyebben eldönthető lehessen, hogy csalt-e a játékos. Például: Focus loss detection (játékos elhagyja-e az oldalt meccs közben - "tab-out"), pontosabb guess time számítás.
    \item \textbf{Új játékmódok:} Olyan játékmódok bevezetése, amelyek elősegítik a tanulást: CO-OP játékmód, ahol több játékos együttes munkával megismer egy adott érettségi témát; csoportos kompetitív játékmódok.
    \item \textbf{Kérdések bővítése:} Olyan közösségi felület, ahol a játékosok kérdéseket javasolhatnak.
    \item \textbf{Mobilapplikáció:} Android, iOS alkalmazás fejlesztése a jobb felhasználói élmény érdekében, például Flutter segítségével.
    \item \textbf{Badge és Achievement rendszer:} Bizonyos mérföldkövek után a játékos Badge-eket kap, amelyek megjelennek a profiljukon, ezzel növelve a motivációt. (A Badge rendszer adatbázis struktúrája megtalálható az MVP-ben, de a implementálásra nem kerül)
    \item \textbf{Távoli jövő - teljes érettségi felkészítő platform:} Monetizált érettségi felkészítő platform, ahol a jelenleg ismertetett funkciókon kívül megtalálható hivatalos tanári segítség és több erőforrás a tanuló felkészülésének elősegítése érdekében.
\end{itemize}